Mechanistic interpretability (MI) analyses of biological foundation models are easy to get wrong: apparent regulatory signals can reflect co-expression, expression level or reference-network structure, and reviewing such analyses is laborious. We present MI-Workbench, an open-source platform that runs critique-and-revision loops on MI analyses with large language model (LLM) agents. Through an inspectable orchestrator, an executor performs and revises an analysis, optionally running its code in a sandbox, while a panel of reviewer lenses (rigour, adversarial, biological plausibility) critiques it. We built a benchmark of 12 flawed analysis write-ups of single-cell foundation models (48 planted flaws) and 6 clean controls, and evaluated 1,296 reviewer calls from three OpenAI model configurations with two judges (flaw-detection agreement $\kappa=0.84$). Three reviewer calls found more flaws than one (90--91\% vs 81--86\%), but three distinct lenses found no more than one lens sampled three times; distinct lenses mainly reduced noise and output tokens. Reviewers caught most statistical, confounding, causal and biological flaws, including the false premise that attention is symmetric, but only about one in five false statements about the analysed model. Word-overlap merging of duplicate critiques failed at an uncalibrated threshold (F1 0.02) but nearly matched LLM adjudication once calibrated (held-out 0.80 vs 0.82). Real revision loops fixed most planted flaws in one revision but never satisfied the panel, so the default is a fixed budget. In an executed case study on Geneformer attention and Tabula Sapiens data, every reported number traced to executed code, and the panel identified at least 25 analysis errors beyond execution failures. Two of three runs agreed with an independent analysis that attention carries no signal about curated TRRUST edges beyond co-expression, rank and hub structure; runs also showed scope creep and hard-coded verdicts. MI-Workbench makes such analyses executable and auditable, but its conclusions require human acceptance.
